%% kapa: exam=1 round=30 year=2019 period=1 form=A subject=economics subject_ko=경제학원론 seq=11 number=51
%% source: Q-Net 감정평가사 기출문제 (cst00309) artlSeq=5205205 file=extracted/1008/5205205/2019년 제30회 감정평가사 1차시험 문제지/2019년 제30회 감정평가사 1차시험 문제지/2019년도 제30회 감정평가사 1차 1교시 A형.hwp
%% conversion: kordoc 4.24.51 -> tools/gen1.py (auto); review=proofread-sonnet; flags=none
\begin{question}{51}
설문을 어떻게 구성하느냐에 따라 다른 응답이 나오는 효과는?
\begin{choices}
  \item 틀짜기효과(framing effect)
  \item 닻내림효과(anchoring effect)
  \item 현상유지편향(status quo bias)
  \item 기정편향(default bias)
  \item 부존효과(endowment effect)
\end{choices}
\end{question}
\answer{1}
